%-*-tex-*-
% Copyright Michael J. Ferguson, INRS-Telecommunications
% All rights reserved. 

% ======== Table of Contents, Figures, and Tables -- Examples ======
% The "list" file is written using \writelistfile and the \shtoc ...
% are used to create the actual table of contents, figures, tables. 
% Extensions to other forms should be obvious. 



%====== general structure for making table of contents, tables, etc.
% ------ table of contents is slightly different because it refers to itself

% file names (defaults) \jobname.toc  --- table of contents
%                       \jobname.fig  --- list of figures
%                       \jobname.tbl  --- list of tables

% to initiate a build of any of these it is only necessary to 
% \openlistfile{<...>} --- This should be only done once 
% when this is done a <..>list flag or switch is also created 

% To print any of these lists it is only necessary to
% use the following format 
%    \begingroup
%     \def\sshtoc#1#2#3{} % omits or specially defines . and any other
%     \preludehead{<name list>} % .... or what you wish
%     \maketoclist % This is for toc files only ... self referential
%     \input \jobname.<extension>
%    \ejectpage % if desired
%    \endgroup

%  a few special macros are defined for toc/list file management (VAX/VMS)

\newif\ifnotdefined % \ifundefined is sensitive if it fails 
\def\checklistopen#1{\notdefinedfalse\ifundefined{if#1list}\notdefinedtrue\relax
                    \fi}
\def\c@k#1{\csname if#1list\endcsname \jobname.#1;-1 \else \jobname.#1 \fi}
\def\newlistfilename#1{\ifnotdefined \jobname.#1 \else \c@k{#1} \fi}
                      % \c@k is necessary because of the way TeX expands 

\newtoks\inputtocfiles
\inputtocfiles = {\checklistopen{toc}\inputwithcheck{\newlistfilename{toc}}}
\def\maketoclist{\begingroup\captionbodyoff\the\inputtocfiles\endgroup}

\newtoks\inputfigfiles
\inputfigfiles = {\checklistopen{fig}\inputwithcheck{\newlistfilename{fig}}}
\def\makefiglist{\begingroup\captionbodyoff\the\inputfigfiles\endgroup}

\newtoks\inputtblfiles
\inputtblfiles = {\checklistopen{tbl}\inputwithcheck{\newlistfilename{tbl}}}
\def\maketbllist{\begingroup\captionbodyoff\the\inputtblfiles\endgroup}


% ======== macro forms for toc, fig, tbl =========
%   Some of these are defined as null to suppress printing of the
%   corresponding entries.
%   #1  Number
%   #2  Text
%   #3  Page number
% ========== toc macro forms ==========
% special forms that allow for multiple line titles ... auto line breaking
% only
%#1--indent  #2 --section number #3--title 
% #4-- pagenumber 
\newdimen\t@ocsize
\newdimen\righttocindent %\righttocindent = 5em

\def\tocfill{\leaders\hbox to .5em{\hss.\hss}\hfill\hskip .5em}

% default font is \rm ... may be changed in the actual \figtoc ...etc
% macros as given below. 

\def\tocbox#1#2#3#4{\line{\rm \t@ocsize=\hsize
                      \setbox2=\hbox {\hskip#1\relax#2\quad}
                      \advance\t@ocsize by -\righttocindent
                      \vbox{\hsize \t@ocsize \normalbaselines
                              \setstrut
                              \hangindent 1\wd2 
                               \hangafter 1 \def\cr{\linebreak}\frenchspacing
                               \veryraggedright \parindent=0pt\leavevmode
                      \box2\relax#3\strut\null\nobreak\rm\tocfill 
                              #4\hskip -\righttocindent\strut}\hss}}
    
% ---------- figure forms -------
\def\figtoc#1#2#3{   \vskip 1ex plus 1pt minus 1pt
                    \tocbox{2em}{#1}{#2}{#3}}
% ---------- table forms -------
\let\tbltoc=\figtoc

\def\prtoc#1#2#3{   \vskip 1ex plus 1pt minus 1pt
                    \tocbox{0em}{\it#1}{\it#2}{\it#3}}

\def\chtoc#1#2#3{   \goodbreak\vskip 2ex plus 1pt minus 1pt
                    \tocbox{0em}{\sheadfont#1}{\sheadfont\def\cr{ }#2}{#3}}
\def\shtoc#1#2#3{   \vskip 1ex plus 1pt minus 1pt
                    \tocbox{2em}{#1}{#2}{#3}}
\def\sshtoc#1#2#3{  \vskip .5ex plus 1pt minus 1pt
                    \tocbox{4em}{#1}{#2}{#3}}
\def\intssshtoc#1#2#3{  \vskip .3ex plus 1pt minus 1pt
                    \tocbox{6em}{#1}{#2}{#3}}
\def\intdssshtoc#1#2#3{  \vskip 1pt plus 1pt minus 1pt
                    \tocbox{8em}{#1}{#2}{#3}}
\def\ssshtoc#1#2#3{}
\def\dssshtoc#1#2#3{}

% the upper two forms are defaulted to be omitted. They are easily included
% with a \let\ssshtoc = \intssshtoc  ... etc . 